\scanpage{91}
\mspage{44}

$\exists$ générateur $g\in T$ (i.e. $g\in T^*$) et un ordre
sur l'ens. $I$, \inserted{$i_1-\cdots-i_{\unsure{\nu}}$} \deleted{\illegible{2}} des conjugués
$\unsure{l_i'|l_{i_\alpha}'}$ des $\unsure{\kappa_{i_\alpha}|k_{i_\alpha}}(g)$, et $g\in\mathbb N$ et des $x_\alpha,y_\alpha\in\pi$
$(1\leq\alpha\leq g)$ tels que l'on ait
\[
\left(\prod_1^g[x_\alpha,y_\alpha]\right)\prod_1^{\unsure{\nu}}l_i=1
\]
et que ceci soit une relation de définition
de $\pi$.

Il y a \inserted{cependant} un grain de sel pour $I=\varnothing$
(\unsure{auquel cas} $T$ ne sert à rien \inserted{\unsure{apparemment|proprement}} (dans la
description groupoïdale de $(X,S)$ \unsure{$\simeq$} \unsure{$(X,\varnothing)(=X)$}))
il faut alors, \deleted{\illegible{1}} \unsure{au lieu} des données
``homomorphismes'', se donner
\[
H^2(\pi,\mathbb Z)\overset{\unsure{\kappa|k}}{\simeq}T
\]
\margin{Il faut introduire \inserted{\unsure{donc}} aussi comme donnée supplémentaire, $\unsure{\kappa|k}$,
dans la définition des groupes à lacets \unsure{l'exclusion des cas} $X_0\simeq S^2$ (\unsure{et alors} $S$ vide) est \illegible{2} particulièrement \illegible{1}.}

Enfin, il faut \deleted{\illegible{1}} \unsure{bien un} \illegible{1}
(\unsure{grain de sel}) exclure le cas $X=S^2$, $S=\varnothing$
(en tant que comp. connexe) --- i.e.
du côté groupes, le cas d'une \inserted{comp. de $\Pi_U$ où}
$\pi=(1)$ \deleted{\illegible{2} $I=\varnothing$} et $I=\varnothing$.
[Si je me rappelle bien, il n'y a pas lieu
d'exclure $(S^2,\mathrm{pt})$ i.e. $X\setminus S=U\simeq\mathbb R^2\simeq\unsure{E^2_{\mathrm e}}$)
\margin{90}

\scanpage{92}

\unsure{un point} i.e. $\pi=0$. Mais sauf dans ces
deux cas, \deleted{les hom. $\unsure{\kappa_i|k_i}:$} (correspondant
\unsure{au cas} $g=0$, $\nu=1$) l'hom. $\unsure{\kappa_i|k_i}:T\longrightarrow\pi$
est injectif. Si $T=\mathbb Z$ la donnée des $\unsure{\kappa_i|k_i}$
équivaut à celle d'éléments $l_i\in\pi_i$, et on
retrouve la définition usuelle des
structures à lacets.

Ceci est explicité dans la thèse de
Yves Ladegaillerie --- ce qui y manque, est
(entre autres) la considération de flèches
entre $(X,S)$ autres que des homéomorphismes
(mod isotopies), p.ex. des applications \deleted{étales}
\deleted{\illegible{3}} $X'\xrightarrow{f}X$ telles que $S'=f^{-1}(S)$
[et que $X'$ soit étale sur $X\setminus S$, \unsure{si on y tient!}],
ce qui se ramène au cas précédent
--- le cas \inserted{plus} général où on suppose seule-
ment que $X'$ est un rev. ramifié
de $X$ (pouvant être ramifié aussi en
dehors de $S'$) et $S'\supset f^{-1}(S)$ (mais $S'$
pouvant être plus grand) demanderait
\margin{91}

\scanpage{93}
\mspage{45}

une étude soigneuse, avec une notion
ad-hoc de \unsure{l'isotopie} ....

Une autre direction importante (notam-
ment pour l'étude \unsure{du|des} cas non orienté,
\unsure{ms} orientable) est l'introduction des
groupes finis d'homéomorphismes, ne
respectant pas nécessairement l'orientation.
Pour traiter le cas du changement d'orientation,
notons que dans la description groupoïdale
$(\Pi_{D^*}\longrightarrow\Pi_U,\unsure{\kappa|k})$ d'un $(X,S)$ $T$-orienté,
le passage à l'orientation opposée s'exprime en
gardant tel quel $\Pi_{D^*}\longrightarrow\Pi_U$, et en
remplaçant $\unsure{\kappa|k}$ par $\unsure{\bar\kappa=\kappa^{-1}|\bar k=k^{-1}}$
\[
\unsure{\bar\kappa|\bar k}(\unsure{\zeta|\xi})=\unsure{\kappa|k}(-\unsure{\zeta|\xi})=\unsure{\kappa|k}(\unsure{\zeta|\xi})^{-1}
\]
(Si \inserted{$T=\mathbb Z$, sur} le système $(\pi,(l_i))$, cela revient à
remplacer les $l_i$ par les $l_i^{-1}$ ...)
\margin{\unsure{Si l'on} (si $I=\varnothing$) $\unsure{\kappa|k}:T\simeq H^2(\pi,\mathbb Z)$
\unsure{rempl.} \illegible{1} $-\unsure{\kappa|k}$.}

Soit donc $\Gamma$ un groupe (a priori
\unsure{non} néc. fini) qui opère sur $(X,S)$, donc
sur $U=X\setminus S$ \inserted{et sur $\Pi_U$, $\Pi_{D^*}$.} On trouve alors des
groupoïdes \inserted{fondamentaux} (mixtes) \deleted{\illegible{1}} par la construction bien connue
\margin{92}

\scanpage{94}

\diagram{\begin{tikzcd}[ampersand replacement=\&]
\Pi_{D^*,\Gamma} \arrow[r] \& \Pi_{U,\Gamma} \arrow[r] \& \Pi_{e,\Gamma}
\end{tikzcd}}
catégorie des $\Gamma$-torseurs
i.e. des objets $1$-connexes
dans $\Gamma$-Ens

correspondant aux foncteurs en sens inverse

\diagram{\begin{tikzcd}[ampersand replacement=\&]
\text{$\Gamma$-rev. \unsure{étales} de $D^*$} \& \text{$\Gamma$-rev. \unsure{étales} de $U$} \arrow[l] \& \text{$\Gamma$-Ens} \arrow[l]
\end{tikzcd}}

Les comp. connexes de $\Pi_{D^*,\Gamma}$ correspondent
aux \uline{orbites de $\Gamma$ dans $S$} \deleted{\illegible{1}} $\simeq\pi_0(\Pi_{D^*})$,
de même pour celles de $\Pi_{U,\Gamma}$. Le \deleted{\unsure{top.}}
$\Pi_{U,\Gamma}$ connexe \inserted{i.e. le topos $\widehat\Pi_{U,\Gamma}$ connexe,}
\unsure{est celui à} $\Gamma$ transitif
sur $\pi_0(U)\simeq\pi_0(X)$ --- quitte à remplacer
$X$ par une composante connexe $X_0$, et $\Gamma$
par le \unsure{s.-s.-groupe|s.-groupe} $\Sigma\subset\Gamma$ qui la stabilise,
on serait ramené pour l'étude du
topos $\widehat\Pi_{U,\Gamma}$ et des morphismes de topos
\diagram{\begin{tikzcd}[ampersand replacement=\&]
\widehat\Pi_{D^*,\Gamma} \arrow[r,dash] \& \widehat\Pi_{U,\Gamma} \arrow[r] \& (\unsure{\widehat\Pi_{e,\Sigma}} \arrow[r] \& )\widehat\Pi_{e,\Gamma}
\end{tikzcd}}
au cas de $(X_0,S_0,\Sigma)$. Mais l'analogue
\margin{93}

\scanpage{95}
\mspage{46}

que j'ai en vue avec le cas ``arithmétique''
rend cette réduction inopportune dans le
cas général (car dans le cas arithmétique,
on ne se borne pas non plus au cas où
$\Sigma=\Gamma$ i.e. $\unsure{\mathcal E}\longrightarrow\Gamma$ surjectif, i.e. \inserted{la comp. connexe}
$X$ \uline{géométrique-}
\uline{ment} connexe sur $K$). Notons \inserted{ici} que
$(X,S,\Gamma)$ se récupère à partir de
$(X_0,S_0,\Gamma_0=\Sigma)$ comme somme \unsure{amalgamée}
\[
X=X_0\unsure{\widehat\times_{\Sigma}|\wedge_{\Sigma}}\Gamma\quad ....
\]

Quand on exprime (pour $\Pi_{\unsure{U|X},\Gamma}$ connexe)
la situation en termes de théorie des groupes,
on trouve donc \deleted{\illegible{1}} un groupe
fondamental \deleted{$\unsure{\mathcal E}_X$} mixte
\[
\unsure{\mathcal E}_{\unsure{U|X},\Gamma}=\pi_1(U,\Gamma)
\]
et un hom.
\[
\unsure{\mathcal E}_{U,\Gamma}\longrightarrow\Gamma
\]
--- (Surjectif ssi. \deleted{$\Gamma$ transitif sur $\pi_0(U)=\unsure{\Gamma/\Sigma}$}
\inserted{$U$ connexe (on a $\Gamma/\Sigma\simeq\pi_0(U)$)})
\inserted{enfin un ens. d'indices}
\deleted{correspondant à \illegible{1} $\pi_0(U)/\Gamma$}
$I$ ($\simeq\unsure{S}/\Gamma\simeq\pi_0(\Pi_{D^*,\Gamma})$) et des
groupes fondamentaux mixtes
\margin{94}

\scanpage{96}

\[
\unsure{\mathcal E}_i\xrightarrow{\sim}\unsure{\pi_1|\Pi_1}(D_i^*,\Gamma)
\simeq\pi_1(D_{i_0}^*,\unsure{\Sigma_{i_0}|\Gamma_{i_0}})
\]
où $D_{i_0}^*$ est une composante connexe du
multidisque troué $D_i^*$ \inserted{(correspondant au choix d'un $\unsure{s_i}\in S$)}, et où $\unsure{\Sigma_{i_0}|\Gamma_{i_0}}\subset\Gamma$ est
son stabilisateur (i.e. le stabilisateur de
$s_i$ dans $\Gamma$), qui est \deleted{\unsure{un gr.}} (si $\Gamma_i$ \inserted{est fini et} opère
fidèlement au voisinage de $s_i$) un groupe cyclique
ou \unsure{dihédral|dihedral} ... On trouve donc, si $X$ est
$T$ orienté, une extension de $\Gamma_{i_0}$ par $T$,
l'hom. $\unsure{\mathcal E}_i\longrightarrow\unsure{\Sigma_{i_0}|\Gamma_{i_0}}\subset\unsure{\Gamma|\Gamma_i}$ étant induit
bien sûr via $\unsure{\mathcal E}_{U,\Gamma}\longrightarrow\unsure{\Gamma}$ et $\unsure{\mathcal E}_i\longrightarrow\unsure{\mathcal E}_{U,\Gamma}$.

Il faut \inserted{encore} lier l'action de $\Gamma$ sur $X$ à
l'orientation de $X$ --- pour ceci on suppose
\unsure{donné} une \unsure{condition}
\[
\chi:\Gamma\longrightarrow\mathbb Z^*=\{\pm1\}
\]
et on exige que pour $g\in\Gamma$, $g_X$ \unsure{conserve}
l'orientation si $\chi(g)=1$, la renverse si $\chi(g)=-1$.
Ceci implique que \deleted{\unsure{l'}}\inserted{l'hom. ou} isomorphisme
\[
\pi_1(\Pi_{D^*,\Gamma}/\Pi_{\unsure{e},\Gamma})
\overset{\unsure{\kappa_\Gamma|k_\Gamma}}{\xleftarrow{\sim}}T
\]
\margin{système local des noyaux des
$\pi_1(\unsure{s|\xi})\longrightarrow\unsure{\pi_1(p(s))|\pi_1(\varphi p(\xi))}$
pour $\unsure{s|\xi}\in\operatorname{Ob}\Pi_{D^*,\Gamma}$.}
\margin{95}

\scanpage{97}
\mspage{47}

$T$ étant considéré comme système local
sur $\Pi_{\Gamma,e}$ \deleted{\illegible{1}} i.e. comme $\Gamma$-groupe, grâce à l'action
de $\Gamma$ via la \unsure{condition} $\chi$. Il revient au
même de dire que \inserted{$\forall i\in I$} l'application
\deleted{La donnée}
\[
\unsure{\kappa_i|k_i}:T\xrightarrow{\sim}L_i=\operatorname{Ker}(\unsure{\mathcal E}_i\longrightarrow\Gamma)
\]
en tant que hom. de $T$ dans $\unsure{\mathcal E}_i$ est
compatible avec l'action de $\unsure{\mathcal E}_i$ (opérant
sur $T$ via $\chi p_i$ ($p_i:\unsure{\mathcal E}_i\longrightarrow\Gamma$)), \unsure{et un}
\illegible{2} par autom. int ...

Si on exclut le cas où $(X_0,\unsure{S_0})$
est $\simeq(S^2,\varnothing)$ ou $(S^2,1\ \mathrm{pt})$ (i.e.
le cas $\pi=0$), les $\unsure{\kappa_i|k_i}:T\longrightarrow\pi$ sont injectifs,
donc aussi les $\unsure{\mathcal E}_i\longrightarrow\unsure{\mathcal E}_{U,\Gamma}$, donc les
$\unsure{\mathcal E}_i$ peuvent être considérés comme des
\uline{sous-groupes} de $\unsure{\mathcal E}_{U,\Gamma}$.

\unsure{\S} Ici il \deleted{\illegible{1}} serait particulièrement
contre-indiqué (\illegible{1} bien \unsure{supposer|suppose} $\Sigma=\Gamma$ i.e.
$X=X_0$ i.e. $X$ connexe) de supposer que
les $\Gamma_i$ soient $=\Gamma$ i.e. que les $s\in S$ sont fixes
par $\Gamma$ !
\margin{96}

\scanpage{98}

Comme le centre de $\pi$ est réduit à $1$
(si on excepte \deleted{\unsure{resp.}} le cas $(X_0,S_0)\simeq(S^2,\text{deux pts})$
i.e. $U_0\simeq\mathbb C^*$ \inserted{--- cas de la couronne ---}), la donnée d'une extension
de $\Sigma$ ($\subset\Gamma$) par $\pi$ revient (à \unsure{isom.} unique
près) à la donnée d'un hom.
\[
\unsure{\Sigma|\Gamma}\longrightarrow\mathrm{Aut\,ext}(\pi)\quad\deleted{\text{structure à lacets}}
\]
\deleted{de $\Sigma$ dans le ``groupe de Teichmüller''
relatif à $\pi$ ($\simeq\pi_1(X_0)$)} --- $\unsure{\mathcal E}$ se reconsti-
tue comme l'image inverse de l'extension
\diagram{\begin{tikzcd}[ampersand replacement=\&]
1 \arrow[r] \& \pi \arrow[r] \& \mathrm{Aut}(\pi) \arrow[r] \& \mathrm{Autext}\,\pi \arrow[r] \& 1
\end{tikzcd}}.

Mais ici les \unsure{automs} extérieurs relatifs
aux $\alpha\in\Sigma$ \deleted{envisagés sont} \uline{respectent la structure}
à lacets \deleted{\illegible{1}} de $\pi$, modulo le
signe $\chi(\alpha)$ --- i.e. pour $g\in\unsure{\mathcal E}$, et $s\in S$
$\exists\ \alpha\in\pi$ \inserted{et $s'\in S$ ($s'$ unique !)} tels que
\[
\deleted{\illegible{1}}\quad
\operatorname{int}(g)\unsure{\kappa_s|k_s}(\unsure{\zeta|\xi})
=\operatorname{int}(\alpha)\unsure{\kappa_{s'}|k_{s'}}(\chi(p(g))\unsure{\zeta|\xi})
\]
$\forall\unsure{\zeta|\xi}\in\unsure{T|\mathbb Z}$, i.e.
\[
\operatorname{int}(\alpha^{-1}g)\unsure{\kappa_s|k_s}(\unsure{\zeta|\xi})
=\deleted{\illegible{1}}\unsure{\kappa_s|\kappa_{s'}}(\chi(p(g))\unsure{\zeta|\xi})
\]
\deleted{ce qui implique que $\alpha^{-1}g$ normalise $L_s=\unsure{\kappa_s|k_s}(T)$}
\margin{97}

\scanpage{99}
\mspage{48}

\deleted{(donc sans doute $\alpha^{-1}g\in\unsure{\mathcal E}_s$ ?), et y
opère \illegible{2}}
Ainsi, l'opération de
$\Sigma$ sur $S_0$ (donc de $\Gamma$ sur $S=S_0\wedge_{\Sigma}\Gamma$)
est connue, par l'opération ext. de $\Sigma\subset\Gamma$ sur
$\pi$ munie de sa structure \deleted{\illegible{1}} à
``$T$-lacets'' (i.e. les hom. extérieurs
$\unsure{\kappa_s|k_s}:T\longrightarrow\pi$) --- donc aussi $I=S_0/\Sigma\simeq S/\Gamma$.

Peut-on reconstituer $\unsure{\mathcal E}_i$ \inserted{($i\in I$)} à partir de la
\deleted{\unsure{si} $i=s\bmod\Sigma$}
structure d'extension ? On voit, \deleted{que}
\deleted{$\unsure{\mathcal E}_{\unsure{i|i_0}}\subset\operatorname{Norm}_{\unsure{\mathcal E}}(L_s)$}
en vertu des choix faits,
que si $i\in I$ (donc $i$ une orbite de
$\Sigma$ dans $S_0$) $\exists\ s\in i$ tel que $\operatorname{Ker}(\unsure{\mathcal E}_i\longrightarrow\Gamma)$
ne soit autre que $L_s=\unsure{\kappa_s|k_s}(T)$, et on a
donc $\unsure{\mathcal E}_i\subset\operatorname{Norm}_{\unsure{\mathcal E}}(L_s)$, Mais
on a
\[
\operatorname{Norm}_{\unsure{\mathcal E}}(L_s)\cap\pi
=\operatorname{Norm}_{\pi}(L_s)=L_s
\]
$\unsure{=|\ne}\unsure{\mathcal E}_i\cap\pi$, et d'autre part l'image de
$\operatorname{Norm}_{\unsure{\mathcal E}}(L_s)$ dans $\Sigma\subset\Gamma$ n'est \deleted{autre que} \inserted{$\unsure{\subset}$}
$\Sigma_s=$ stabilisateur \deleted{de} \inserted{$s$ dans} $\Sigma$
$=$ Image de $\unsure{\mathcal E}_i$ dans $\Sigma$,
\unsure{on} \unsure{en} résumé
\[
\unsure{\mathcal E}_i\unsure{\equiv|=}\operatorname{Norm}_{\unsure{\mathcal E}}(L_s)
\]
\margin{98}

\scanpage{100}

Inversement, la donnée de $\unsure{\mathcal E}\longrightarrow\Gamma$
et des $\unsure{\mathcal E}_i$, $\unsure{\kappa_i|k_i}:T\longrightarrow\unsure{\mathcal E}_i$ redonne la
structure à lacets de $\pi$, en prenant
les $\unsure{\kappa_i|k_i}$ et tous leurs conjugués exté-
rieurs \inserted{distincts} par les $\alpha\in\unsure{\Sigma|\mathcal E}$ (modifiés par $\chi(\alpha)$ ...).
Donc la \inserted{donnée de la} situation $\unsure{\mathcal E}\xrightarrow{p}\Gamma$, $\unsure{\mathcal E}_i\subset\unsure{\mathcal E}$
\deleted{équivalent} (famille de ss-groupes, chacun
\unsure{défin.} mod conjugaison \uline{dans $\unsure{\mathcal E}$}) équivaut
à la donnée de

a) $\pi=\operatorname{Ker}p$, avec sa structure à \inserted{$T$-}lacets
(\deleted{famille finie de \unsure{classes de conj.}}\inserted{ens. fini d'hom. extérieurs de}
$T$ dans $\pi$)

b) Un \unsure{sous-}\deleted{groupe} groupe $\unsure{\Sigma|\Gamma}\subset\Gamma$, et
un hom.
\[
\Sigma\longrightarrow\text{Teichmüller étendu de }\pi
\]
(autom. extérieurs de $\pi$,
respectant la structure à lacets
modulo signe
compatible avec la \unsure{condition} $\chi|\Sigma$ et
la \unsure{condition} ``signe'' sur Teichmüller étendu.
En fait, $(X,S,\Gamma)$, ou $(U,\Gamma)$, ne définit
$\pi$ que comme groupe extérieur \inserted{à $T$-lacets} sur lequel
$\Gamma$ opère de façon compatible avec $\chi$.
\margin{99}
